This book has one idea and one dedication.

The idea is that a quantum fluid can be understood through \emph{three instruments used together}. The first is the experiment: a measurement made at the bench, in a cryostat, on a neutron beam line, in a vacuum chamber of cold atoms. The second is a proof: a statement about the mathematics of the model, checked by a computer so that it cannot be wrong in the way an arithmetic slip or a misprinted coefficient can be wrong. The third is a simulation: a solver that integrates the equations of the model and lets us watch what the proof only asserts. Each instrument is blind where the others see. A proof says nothing about helium; an experiment does not tell you whether a printed formula is right; a simulation cannot distinguish a theorem from a coincidence. Used together, and with their limits stated, they give something firmer than any of them alone.

The dedication is to Henri Godfrin. His group at the Institut N\'eel in Grenoble does the hardest kind of physics: it cools matter to microkelvin temperatures, puts a single atomic layer of helium-3, or a bath of helium-4, in a neutron beam, and reads the excitation spectrum of a quantum fluid directly. The measurements of that programme --- the roton-like collective mode of a two-dimensional Fermi liquid, and the dispersion relation of superfluid helium-4 together with its thermodynamic consequences --- are the experimental thread of this book. Where the Lean library formalizes a statement that appears in his papers, the book says exactly which statement and what the proof does and does not show. Nothing here speaks for him: he has not reviewed this book, the proofs are of mathematics and not of his data, and every error that remains is ours.

The Lean~4 library is \texttt{QuantumFluids}, a collection of machine-checked theorems on the mathematical structure of quantum fluids; the solver is rusty-SUNDIALS, a Rust implementation of the SUNDIALS suite with a quantum-fluids module \texttt{qf-pgpe}. Both belong to a research programme that is run under strict discipline: predictions written before data are read, failures reported, claims retracted when they fail. That discipline shapes the book. Every chapter contains a box that says what is established and what is not, and the last chapter tells, plainly, what went wrong in the programme and how it was found. A tribute to an experimentalist should be written in the register of experimentalists: with the error bars in view.

The book is organized in five parts and fifteen chapters that can be read in order or separately. Part~I (\cref{ch01,ch02}) introduces the three instruments and the measurements that motivate them, and is a primer on Lean and on the solver. Part~II, on Bose superfluids, goes from Bose--Einstein condensation (\cref{ch11}) through circulation and vortices, the specific heat of helium-4 and the dispersion relation to two-fluid hydrodynamics and second sound (\cref{ch12}). Part~III follows the vortices: the Kosterlitz--Thouless flow, friction, Onsager's giant vortex clusters, the flow past an obstacle and quantum turbulence (\cref{ch06,ch07,ch13,ch09,ch15}). Part~IV turns to Fermi fluids: the two-dimensional Fermi liquid with its zero sound and roton mode (\cref{ch08}), and pairing and the superfluidity of helium-3 (\cref{ch14}). Part~V (\cref{ch10}) is about trust. Each chapter takes its theme with a measurement, a theorem and a computation, and ends with three exercises whose worked solutions are in Appendix~C. Appendix~A is the complete catalogue of the Lean library, Appendix~B tells how to rebuild every figure and rerun the audit, and Appendix~D collects the constants, units and notation. Every figure in the book is produced by a script that is part of the source, and every number quoted in the text comes from a file or a computation recorded alongside it.

\paragraph{How the book was made.} The chapters were written separately, each from the same briefing, and combined at the end; the text, the Lean code and the figures were drafted with the assistance of an AI system and then checked: the Lean files compile against the pinned Mathlib, the solver code was run, and the sources of every number are listed in the per-chapter reports kept with the source. The second edition adds five chapters, the solutions of all the exercises and the appendix of constants and notation, made and checked in the same way; the editor's corrections are dated in the briefing files kept with the source. Corrections are welcome.

\begin{flushright}\itshape X.~C.\\ October 2026\end{flushright}
